\documentclass[12pt, a4paper]{article}

% ─── Codificación y tipografía ───────────────────────────────────────────────
\usepackage[utf8]{inputenc}
\usepackage[T1]{fontenc}
\usepackage{babel}
\babelprovide[main, import]{spanish}
%\usepackage{lmodern}

% ─── Matemática ──────────────────────────────────────────────────────────────
\usepackage{amsmath, amssymb, amsthm}
\usepackage{bm}          % vectores en negrita

% ─── Gráficos y colores ──────────────────────────────────────────────────────
\usepackage{xcolor}
\usepackage{graphicx}
\usepackage{float}

% ─── Tablas ──────────────────────────────────────────────────────────────────
\usepackage{booktabs}
\usepackage{array}

% ─── Hipervínculos ───────────────────────────────────────────────────────────
\usepackage[colorlinks=true, linkcolor=unsblue, urlcolor=unsblue, citecolor=unsblue]{hyperref}

% ─── Márgenes ────────────────────────────────────────────────────────────────
\usepackage[top=2.5cm, bottom=2.5cm, left=3cm, right=3cm, headheight=15pt]{geometry}

% ─── Encabezado y pie ────────────────────────────────────────────────────────
\usepackage{fancyhdr}
\pagestyle{fancy}
\fancyhf{}
\renewcommand{\headrulewidth}{0.4pt}
\fancyhead[L]{\small\textcolor{unsblue}{\textbf{MAD1} --- Métodos para el Análisis de Datos I}}
\fancyhead[R]{\small\textcolor{unsblue}{Unidad 4: Aprendizaje Supervisado}}
\fancyfoot[C]{\thepage}

% ─── Color institucional UNS ─────────────────────────────────────────────────
\definecolor{unsblue}{RGB}{30,60,120}
\definecolor{lightblue}{RGB}{220,230,245}
\definecolor{alertred}{RGB}{180,30,30}
\definecolor{defgreen}{RGB}{20,100,50}

% ─── Entornos teóricos ───────────────────────────────────────────────────────
\theoremstyle{definition}
\newtheorem{definicion}{Definición}[section]
\newtheorem{ejemplo}{Ejemplo}[section]

\theoremstyle{plain}
\newtheorem{proposicion}{Proposición}[section]

\theoremstyle{remark}
\newtheorem{observacion}{Observación}[section]
\newtheorem{nota}{Nota}[section]

% ─── Caja de concepto clave ──────────────────────────────────────────────────
\usepackage{tcolorbox}
\tcbuselibrary{skins, breakable}

\newtcolorbox{conceptoclave}[1][]{
  colback=lightblue,
  colframe=unsblue,
  fonttitle=\bfseries\color{unsblue},
  title={Concepto clave},
  breakable,
  #1
}

\newtcolorbox{alerta}[1][]{
  colback=red!5,
  colframe=alertred,
  fonttitle=\bfseries\color{alertred},
  title={Atención},
  breakable,
  #1
}

% ─── Numeración de secciones ─────────────────────────────────────────────────
\setcounter{secnumdepth}{3}

% ─── Interlineado ────────────────────────────────────────────────────────────
\usepackage{setspace}
\setstretch{1.15}

% ─── Comandos propios ────────────────────────────────────────────────────────
\newcommand{\R}{\mathbb{R}}
\newcommand{\E}{\mathbb{E}}
\newcommand{\bx}{\bm{x}}
\newcommand{\bw}{\bm{w}}
\newcommand{\by}{\bm{y}}
\newcommand{\bX}{\bm{X}}
\newcommand{\yhat}{\hat{y}}
\newcommand{\sign}{\operatorname{sign}}

% ─────────────────────────────────────────────────────────────────────────────
\begin{document}
% ─────────────────────────────────────────────────────────────────────────────

% ── Portada ──────────────────────────────────────────────────────────────────
\begin{titlepage}
  \begin{center}
    \vspace*{1.5cm}
    {\Large \textbf{\textcolor{unsblue}{Universidad Nacional del Sur}}}\\[0.3cm]
    {\large Departamento de Matemática}\\[1.5cm]

    \rule{\textwidth}{1.5pt}\\[0.4cm]
    {\LARGE \textbf{\textcolor{unsblue}{Métodos para el Análisis de Datos I}}}\\[0.4cm]
    {\large Código 8142 \, | \, Licenciatura en Matemática Aplicada}\\[0.3cm]
    \rule{\textwidth}{1.5pt}\\[1.5cm]

    {\Huge \textbf{Unidad 4}}\\[0.5cm]
    {\LARGE \textbf{\textcolor{unsblue}{Aprendizaje Supervisado}}}\\[2.5cm]

    {\large Dr.\ José Bavio}\\[0.3cm]
    {\large 2026}
  \end{center}
\end{titlepage}

% ── Tabla de contenidos ──────────────────────────────────────────────────────
\tableofcontents
\newpage

% ─────────────────────────────────────────────────────────────────────────────
\section{Introducción al Aprendizaje Supervisado}
% ─────────────────────────────────────────────────────────────────────────────

El \textbf{aprendizaje automático} (\textit{machine learning}, ML) es una rama de la inteligencia artificial que estudia algoritmos capaces de aprender patrones a partir de datos sin ser programados explícitamente para cada tarea.

Dentro del ML se distinguen tres grandes paradigmas:

\begin{itemize}
  \item \textbf{Aprendizaje supervisado:} el modelo aprende a partir de pares entrada--salida etiquetados $(x_i, y_i)$.
  \item \textbf{Aprendizaje no supervisado:} no hay etiquetas; el objetivo es descubrir estructura en los datos (Unidad 6).
  \item \textbf{Aprendizaje por refuerzo:} un agente aprende tomando decisiones y recibiendo recompensas.
\end{itemize}

\begin{definicion}[Aprendizaje Supervisado]
Dado un conjunto de entrenamiento $\mathcal{D} = \{(\bx_i, y_i)\}_{i=1}^{n}$ donde $\bx_i \in \R^p$ es el vector de \emph{features} e $y_i$ es la variable respuesta, el objetivo es aprender una función $f: \R^p \to \mathcal{Y}$ que generalice bien a datos nuevos.
\end{definicion}

Según el tipo de variable respuesta:
\begin{itemize}
  \item Si $y_i \in \R$ \quad $\Rightarrow$ \textbf{regresión} (cubierta en Unidad 3).
  \item Si $y_i \in \{1, \ldots, K\}$ \quad $\Rightarrow$ \textbf{clasificación} (tema central de esta unidad).
\end{itemize}

\subsection{El paradigma de generalización}

Aprender de los datos de entrenamiento no es el objetivo final: lo que importa es que el modelo \emph{generalice}, es decir, que funcione bien sobre datos que no vio durante el entrenamiento.

\begin{conceptoclave}
  El \textbf{error de generalización} (o error de prueba) es la pérdida esperada sobre nuevas observaciones:
  \[
    \text{Err} = \E\bigl[\ell(y, f(\bx))\bigr]
  \]
  donde $\ell$ es la función de pérdida (por ejemplo, pérdida 0-1 para clasificación). El objetivo del aprendizaje supervisado es minimizar $\text{Err}$, no el error de entrenamiento.
\end{conceptoclave}

\subsection{Bias-variance tradeoff}

Para modelos de regresión se puede demostrar que el error cuadrático medio de predicción descompone como:
\[
  \text{ECM}\bigl[f(\bx_0)\bigr]
  = \underbrace{\text{Var}\bigl[\hat{f}(\bx_0)\bigr]}_{\text{varianza}}
  + \underbrace{\bigl[\text{Bias}(\hat{f}(\bx_0))\bigr]^2}_{\text{sesgo}^2}
  + \underbrace{\sigma^2_\varepsilon}_{\text{irreducible}}
\]

La misma tensión existe en clasificación:
\begin{itemize}
  \item Modelos \textbf{demasiado simples} tienen alto sesgo (\textit{underfitting}).
  \item Modelos \textbf{demasiado complejos} tienen alta varianza (\textit{overfitting}).
\end{itemize}

\begin{figure}[H]
  \centering
  % Diagrama esquemático: reemplazar con figura real si se dispone
  \fbox{\parbox{0.7\textwidth}{\centering\vspace{1cm}
    \textit{[Insertar curva de bias-variance vs.\ complejidad del modelo]}
  \vspace{1cm}}}
  \caption{Compromiso sesgo-varianza en función de la complejidad del modelo.}
\end{figure}

\subsection{Partición de datos y validación cruzada}

Para estimar el error de generalización se divide el conjunto de datos:

\begin{itemize}
  \item \textbf{Entrenamiento} (\textit{training set}): ajuste de los parámetros del modelo.
  \item \textbf{Validación} (\textit{validation set}): selección de hiperparámetros.
  \item \textbf{Prueba} (\textit{test set}): evaluación final; se toca una \emph{única vez}.
\end{itemize}

La \textbf{validación cruzada $k$-fold} divide el conjunto de entrenamiento en $k$ partes (\textit{folds}) iguales; en cada iteración se entrena con $k-1$ folds y se evalúa en el restante. El error de validación cruzada es el promedio de los $k$ errores:
\[
  \text{CV}_k = \frac{1}{k} \sum_{j=1}^{k} \text{Err}_j
\]

\begin{observacion}
  La validación cruzada \emph{estratificada} (stratified $k$-fold) preserva la proporción de clases en cada fold, lo que es fundamental con clases desbalanceadas.
\end{observacion}

% ─────────────────────────────────────────────────────────────────────────────
\section{Árboles de Decisión}
% ─────────────────────────────────────────────────────────────────────────────

\subsection{Concepto y estructura}

Un \textbf{árbol de decisión} es un modelo que particiona el espacio de features mediante reglas de la forma $x_j \leq t$ (\textit{splits}) organizadas en una estructura jerárquica de árbol.

\begin{definicion}[Árbol de clasificación]
  Un árbol de clasificación $T$ es una estructura formada por:
  \begin{itemize}
    \item \textbf{Nodo raíz:} contiene todo el conjunto de entrenamiento.
    \item \textbf{Nodos internos:} cada uno define un split $\{x_j \leq t\}$ que divide los datos en dos subconjuntos (subárboles izquierdo y derecho).
    \item \textbf{Hojas:} regiones terminales $R_m$ que asignan una clase $\hat{y} = \arg\max_k \hat{p}_{mk}$, donde $\hat{p}_{mk}$ es la proporción de observaciones de clase $k$ en la hoja $m$.
  \end{itemize}
\end{definicion}

\subsection{Criterios de partición}

Para elegir el mejor split $(\bm{j}^*, t^*)$ se minimiza la impureza de los nodos hijos. Los criterios más usados son:

\subsubsection{Índice de Gini}

\[
  G(p) = \sum_{k=1}^{K} p_k(1 - p_k) = 1 - \sum_{k=1}^{K} p_k^2
\]

donde $p_k$ es la proporción de la clase $k$ en el nodo. $G = 0$ indica nodo puro.

\subsubsection{Entropía (ganancia de información)}

\[
  H(p) = -\sum_{k=1}^{K} p_k \log_2(p_k)
\]

La \textbf{ganancia de información} de un split es la reducción en entropía:
\[
  \Delta H = H(\text{padre}) - \frac{n_L}{n} H(\text{hijo}_L) - \frac{n_R}{n} H(\text{hijo}_R)
\]

\begin{observacion}
  Gini y entropía producen resultados muy similares en la práctica. Gini es computacionalmente más eficiente al evitar el cálculo de logaritmos.
\end{observacion}

\subsection{Construcción: algoritmo CART}

El algoritmo \textbf{CART} (\textit{Classification and Regression Trees}) construye el árbol de manera \textit{greedy} (voraz): en cada nodo busca el split que maximiza la reducción de impureza, sin garantía de optimalidad global.

\begin{enumerate}
  \item Iniciar con el nodo raíz conteniendo todos los datos.
  \item Para cada feature $j$ y umbral $t$, calcular la impureza ponderada del split.
  \item Elegir el split $(j^*, t^*)$ que minimiza la impureza total.
  \item Repetir recursivamente en cada subárbol.
  \item Detener cuando se cumple un criterio de parada (profundidad máxima, mínimo de muestras por hoja, etc.).
\end{enumerate}

\subsection{Poda}

Árboles crecidos sin restricción tienen alta varianza. La \textbf{poda por complejidad de costo} (\textit{cost-complexity pruning}) busca el subárbol $T \subset T_{\max}$ que minimiza:
\[
  C_\alpha(T) = \sum_{m=1}^{|T|} \sum_{i \in R_m} \mathbf{1}(y_i \neq \hat{y}_m) + \alpha |T|
\]
donde $|T|$ es el número de hojas y $\alpha \geq 0$ es el parámetro de regularización (elegido por validación cruzada).

\subsection{Ventajas y limitaciones}

\begin{center}
\begin{tabular}{p{0.45\textwidth} p{0.45\textwidth}}
  \toprule
  \textbf{Ventajas} & \textbf{Limitaciones} \\
  \midrule
  Interpretabilidad directa & Alta varianza: pequeños cambios en los datos producen árboles muy distintos \\
  Maneja features categóricos y numéricos & Fronteras de decisión solo axiales \\
  No requiere normalización & Tendencia a sobreajuste si no se poda \\
  Captura interacciones entre variables & Sesgo en features con muchas categorías \\
  \bottomrule
\end{tabular}
\end{center}

% ─────────────────────────────────────────────────────────────────────────────
\section{Random Forest}
% ─────────────────────────────────────────────────────────────────────────────

\subsection{Bagging}

El \textbf{bagging} (\textit{bootstrap aggregating}, Breiman 1996) reduce la varianza promediando múltiples modelos inestables entrenados sobre muestras bootstrap del conjunto de entrenamiento.

Dado un conjunto $\mathcal{D}$ de $n$ observaciones:
\begin{enumerate}
  \item Para $b = 1, \ldots, B$: extraer una muestra bootstrap $\mathcal{D}^{(b)}$ de tamaño $n$ con reemplazo y entrenar un árbol $T^{(b)}$.
  \item Predicción por mayoría de votos:
  \[
    \hat{y}(\bx) = \arg\max_k \sum_{b=1}^{B} \mathbf{1}\bigl(\hat{y}^{(b)}(\bx) = k\bigr)
  \]
\end{enumerate}

\subsection{Random Forest: decorrelación de árboles}

\begin{definicion}[Random Forest]
  \textbf{Random Forest} (Breiman, 2001) extiende el bagging añadiendo \emph{aleatoriedad en la selección de features}: en cada split de cada árbol, se elige el mejor split solo entre un subconjunto aleatorio de $m$ features (en lugar de los $p$ totales).
\end{definicion}

El parámetro $m$ controla el grado de decorrelación:
\begin{itemize}
  \item Valor por defecto para clasificación: $m = \lfloor\sqrt{p}\rfloor$
  \item Valor por defecto para regresión: $m = \lfloor p/3 \rfloor$
\end{itemize}

La decorrelación es esencial: si los árboles fueran muy correlacionados, el promedio no reduciría la varianza.

\begin{proposicion}[Reducción de varianza por promediado]
  Sea $\rho$ la correlación entre dos árboles cualesquiera y $\sigma^2$ la varianza de cada uno. La varianza del promedio de $B$ árboles es:
  \[
    \rho \sigma^2 + \frac{1-\rho}{B}\sigma^2
  \]
  Cuando $B \to \infty$, la varianza tiende a $\rho\sigma^2$. Reducir $\rho$ (mediante la aleatoriedad en features) es lo que diferencia RF del bagging puro.
\end{proposicion}

\subsection{Estimación out-of-bag (OOB)}

Cada árbol $T^{(b)}$ se entrena sobre $\mathcal{D}^{(b)}$. Las observaciones \emph{no incluidas} en $\mathcal{D}^{(b)}$ (aproximadamente $36.8\%$ de $\mathcal{D}$) constituyen la muestra \textbf{out-of-bag} (OOB). El error OOB es un estimador casi insesgado del error de generalización y puede usarse en lugar de validación cruzada.

\subsection{Importancia de variables}

Random Forest provee dos medidas de importancia:

\subsubsection{Importancia por impureza (MDI)}

La importancia de la variable $j$ es la reducción total de impureza (Gini o entropía) generada por los splits en $j$, promediada sobre todos los árboles:
\[
  \text{Imp}(j) = \frac{1}{B} \sum_{b=1}^{B} \sum_{\substack{t \in T^{(b)} \\ \text{split en }j}} \Delta G_t
\]

\subsubsection{Importancia por permutación (MDA)}

Para cada árbol $T^{(b)}$ y cada variable $j$:
\begin{enumerate}
  \item Calcular el error OOB base.
  \item Permutar aleatoriamente los valores de $x_j$ en el conjunto OOB.
  \item Recalcular el error OOB con la permutación.
  \item La importancia es el aumento promedio del error.
\end{enumerate}

\begin{alerta}
  MDI sobreestima la importancia de variables con muchas categorías o alta cardinalidad. Para datos con features de muy distinta naturaleza, se recomienda la importancia por permutación (MDA).
\end{alerta}

\subsection{Hiperparámetros principales}

\begin{center}
\begin{tabular}{lll}
  \toprule
  \textbf{Hiperparámetro} & \textbf{Efecto} & \textbf{Valor típico} \\
  \midrule
  $n\_estimators$ & Número de árboles & 100--500 \\
  $max\_features$ & Features considerados por split & $\sqrt{p}$ (clasif.) \\
  $max\_depth$ & Profundidad máxima del árbol & None (sin límite) \\
  $min\_samples\_leaf$ & Mínimo de muestras por hoja & 1--5 \\
  $bootstrap$ & Usar muestras bootstrap & True \\
  \bottomrule
\end{tabular}
\end{center}

% ─────────────────────────────────────────────────────────────────────────────
\section{Support Vector Machines (SVM)}
% ─────────────────────────────────────────────────────────────────────────────

\subsection{Clasificador de margen máximo}

Para un problema de clasificación binaria con clases $y \in \{-1, +1\}$, se busca el hiperplano
\[
  \mathcal{H} = \{\bx : \bw^\top \bx + b = 0\}
\]
que separa las dos clases con el \textbf{máximo margen}. El margen es la distancia entre el hiperplano y las observaciones más cercanas de cada clase (\textit{vectores de soporte}).

\subsection{Problema de optimización (caso separable)}

\begin{definicion}[SVM de margen duro]
  Minimizar $\frac{1}{2}\|\bw\|^2$ sujeto a
  \[
    y_i(\bw^\top \bx_i + b) \geq 1 \quad \forall\, i = 1, \ldots, n
  \]
\end{definicion}

El margen geométrico es $2 / \|\bw\|$, de modo que minimizar $\|\bw\|^2$ equivale a maximizarlo.

\subsection{SVM de margen blando (soft margin)}

Cuando los datos no son linealmente separables, se introducen \textbf{variables de holgura} $\xi_i \geq 0$:

\[
  \min_{\bw, b, \bm{\xi}}\; \frac{1}{2}\|\bw\|^2 + C\sum_{i=1}^{n}\xi_i
  \quad\text{sujeto a}\quad
  y_i(\bw^\top \bx_i + b) \geq 1 - \xi_i,\quad \xi_i \geq 0
\]

El parámetro $C > 0$ controla el \textbf{compromiso sesgo-varianza}:
\begin{itemize}
  \item $C$ grande: margen más pequeño, menos errores de entrenamiento, más riesgo de sobreajuste.
  \item $C$ pequeño: margen más grande, mayor tolerancia a errores, más regularización.
\end{itemize}

\subsection{Truco del kernel}

La formulación dual del problema SVM depende de los datos solo a través de productos internos $\langle \bx_i, \bx_j \rangle$. Reemplazando el producto interno por una \textbf{función kernel} $K(\bx_i, \bx_j) = \langle \phi(\bx_i), \phi(\bx_j) \rangle$ se proyectan implícitamente los datos a un espacio de mayor dimensión sin calcularlo explícitamente.

\begin{center}
\begin{tabular}{lll}
  \toprule
  \textbf{Kernel} & \textbf{Fórmula} & \textbf{Parámetro(s)} \\
  \midrule
  Lineal & $K(\bx, \bx') = \bx^\top \bx'$ & --- \\
  Polinomial & $K(\bx, \bx') = (\gamma\bx^\top\bx' + r)^d$ & $\gamma, r, d$ \\
  RBF (Gaussiano) & $K(\bx, \bx') = \exp(-\gamma\|\bx-\bx'\|^2)$ & $\gamma$ \\
  Sigmoide & $K(\bx, \bx') = \tanh(\gamma\bx^\top\bx' + r)$ & $\gamma, r$ \\
  \bottomrule
\end{tabular}
\end{center}

\begin{conceptoclave}
  El kernel RBF con $\gamma$ pequeño produce fronteras suaves (mayor regularización); con $\gamma$ grande produce fronteras muy ajustadas al entrenamiento (riesgo de sobreajuste).
\end{conceptoclave}

\subsection{Extensión a clasificación multiclase}

Las SVM son binarias por definición. Las estrategias usuales para $K > 2$ clases son:
\begin{itemize}
  \item \textbf{Uno contra todos (OvR):} se entrena un clasificador por clase.
  \item \textbf{Uno contra uno (OvO):} se entrena un clasificador para cada par de clases; la predicción final es por votación.
\end{itemize}

% ─────────────────────────────────────────────────────────────────────────────
\section{K-Nearest Neighbors (KNN)}
% ─────────────────────────────────────────────────────────────────────────────

\subsection{Definición}

\begin{definicion}[$k$-NN clasificador]
  Dado un punto de consulta $\bx_0$, el clasificador $k$-NN asigna la clase mayoritaria entre sus $k$ vecinos más cercanos en el conjunto de entrenamiento:
  \[
    \hat{y}(\bx_0) = \arg\max_c \sum_{i \in \mathcal{N}_k(\bx_0)} \mathbf{1}(y_i = c)
  \]
  donde $\mathcal{N}_k(\bx_0)$ es el conjunto de los $k$ índices con menor distancia a $\bx_0$.
\end{definicion}

La distancia estándar es la euclídea; alternativas incluyen Manhattan, Minkowski ($L^p$) y distancia de Mahalanobis.

\subsection{Elección de $k$}

\begin{itemize}
  \item $k = 1$: frontera de decisión muy irregular, alta varianza.
  \item $k$ grande: frontera suave, alto sesgo.
  \item Se elige mediante validación cruzada; típicamente se prueban valores impares para evitar empates en clasificación binaria.
\end{itemize}

\subsection{La maldición de la dimensionalidad}

En espacios de alta dimensión ($p \gg 1$), las distancias entre puntos tienden a homogeneizarse: la diferencia entre el vecino más cercano y el más lejano se hace proporcionalmente pequeña. Esto deteriora drásticamente el rendimiento de $k$-NN.

\begin{alerta}
  $k$-NN \textbf{requiere normalización o estandarización} previa de los features, ya que es sensible a la escala de las variables.
\end{alerta}

% ─────────────────────────────────────────────────────────────────────────────
\section{Redes Neuronales Artificiales}
% ─────────────────────────────────────────────────────────────────────────────

\subsection{La neurona artificial (perceptrón)}

Una \textbf{neurona artificial} toma $p$ entradas $x_1, \ldots, x_p$, las combina linealmente y aplica una función de activación $\sigma$:
\[
  a = \sigma\!\left(\sum_{j=1}^{p} w_j x_j + b\right) = \sigma(\bw^\top \bx + b)
\]

\subsection{Perceptrón Multicapa (MLP)}

\begin{definicion}[Perceptrón Multicapa]
  Un \textbf{MLP} (\textit{Multilayer Perceptron}) es una red neuronal con capas totalmente conectadas (\textit{fully connected}):
  \begin{itemize}
    \item Una \textbf{capa de entrada} de dimensión $p$.
    \item Una o más \textbf{capas ocultas} de neuronas con activación no lineal.
    \item Una \textbf{capa de salida} que produce la predicción final.
  \end{itemize}
\end{definicion}

La salida de la capa $l$ es:
\[
  \bm{h}^{(l)} = \sigma^{(l)}\!\left(\bm{W}^{(l)} \bm{h}^{(l-1)} + \bm{b}^{(l)}\right)
\]
donde $\bm{W}^{(l)}$ y $\bm{b}^{(l)}$ son la matriz de pesos y el vector de sesgos de la capa $l$.

\subsection{Funciones de activación}

\begin{center}
\begin{tabular}{llll}
  \toprule
  \textbf{Nombre} & \textbf{Fórmula} & \textbf{Rango} & \textbf{Uso típico} \\
  \midrule
  Sigmoide & $\sigma(z) = 1/(1+e^{-z})$ & $(0,1)$ & Salida binaria \\
  Tanh & $\tanh(z)$ & $(-1,1)$ & Capas ocultas \\
  ReLU & $\max(0, z)$ & $[0,\infty)$ & Capas ocultas (estándar) \\
  Leaky ReLU & $\max(\alpha z, z)$, $\alpha \ll 1$ & $\R$ & Evita neuronas muertas \\
  Softmax & $e^{z_k}/\sum_j e^{z_j}$ & $(0,1)^K$ & Salida multiclase \\
  \bottomrule
\end{tabular}
\end{center}

\subsection{Función de pérdida}

Para clasificación multiclase se usa la \textbf{entropía cruzada categórica}:
\[
  \mathcal{L} = -\frac{1}{n}\sum_{i=1}^{n}\sum_{k=1}^{K} y_{ik} \log \hat{p}_{ik}
\]
donde $y_{ik} \in \{0,1\}$ es el indicador de clase (codificación \textit{one-hot}) y $\hat{p}_{ik}$ es la probabilidad predicha para la clase $k$.

\subsection{Entrenamiento: descenso por gradiente y backpropagation}

Los parámetros $\bm{\theta} = \{\bm{W}^{(l)}, \bm{b}^{(l)}\}$ se ajustan minimizando $\mathcal{L}$ mediante \textbf{descenso por gradiente estocástico} (SGD):
\[
  \bm{\theta} \leftarrow \bm{\theta} - \eta \nabla_{\bm{\theta}} \mathcal{L}
\]
donde $\eta$ es la \textbf{tasa de aprendizaje} (\textit{learning rate}).

El \textbf{algoritmo de backpropagation} calcula eficientemente $\nabla_{\bm{\theta}} \mathcal{L}$ aplicando la regla de la cadena capa por capa (de la salida hacia la entrada).

\begin{conceptoclave}
  Backpropagation no es un algoritmo de optimización por sí mismo: es un método eficiente para calcular el gradiente. El algoritmo de optimización que usa ese gradiente puede ser SGD, Adam, RMSProp, etc.
\end{conceptoclave}

\subsection{Variantes de SGD}

\begin{itemize}
  \item \textbf{Mini-batch SGD:} en cada paso se usa un subconjunto (\textit{batch}) de $m$ observaciones ($m$ típicamente entre 32 y 256).
  \item \textbf{Momentum:} acumula gradientes previos para acelerar la convergencia.
  \item \textbf{Adam:} combina momentum con adaptación individual de la tasa de aprendizaje por parámetro; es el optimizador más usado en la práctica.
\end{itemize}

\subsection{Regularización en redes neuronales}

\begin{itemize}
  \item \textbf{Dropout:} durante el entrenamiento, cada neurona se desactiva con probabilidad $p_{\text{drop}}$, forzando redundancia en la representación.
  \item \textbf{Regularización $L_2$ (weight decay):} añade $\lambda\|\bm{\theta}\|^2$ a la pérdida.
  \item \textbf{Batch normalization:} normaliza las activaciones de cada capa, acelerando el entrenamiento.
  \item \textbf{Early stopping:} detiene el entrenamiento cuando el error de validación deja de decrecer.
\end{itemize}

% ─────────────────────────────────────────────────────────────────────────────
\section{Métricas de Evaluación para Clasificación}
% ─────────────────────────────────────────────────────────────────────────────

\subsection{Matriz de confusión}

Para clasificación binaria (positivo/negativo), la \textbf{matriz de confusión} organiza las predicciones del modelo:

\begin{center}
\begin{tabular}{cc|cc}
  & & \multicolumn{2}{c}{\textbf{Predicho}} \\
  & & Positivo & Negativo \\
  \hline
  \multirow{2}{*}{\textbf{Real}} & Positivo & VP & FN \\
  & Negativo & FP & VN \\
\end{tabular}
\end{center}

donde VP = verdaderos positivos, VN = verdaderos negativos, FP = falsos positivos (error tipo I), FN = falsos negativos (error tipo II).

\subsection{Métricas derivadas}

\begin{align}
  \text{Exactitud (Accuracy)} &= \frac{\text{VP} + \text{VN}}{\text{VP}+\text{FP}+\text{FN}+\text{VN}} \\[6pt]
  \text{Precisión (Precision)} &= \frac{\text{VP}}{\text{VP}+\text{FP}} \\[6pt]
  \text{Sensibilidad (Recall, TPR)} &= \frac{\text{VP}}{\text{VP}+\text{FN}} \\[6pt]
  \text{Especificidad (TNR)} &= \frac{\text{VN}}{\text{VN}+\text{FP}} \\[6pt]
  F_1 &= 2 \cdot \frac{\text{Precisión} \times \text{Recall}}{\text{Precisión} + \text{Recall}}
\end{align}

\begin{alerta}
  La exactitud es una métrica \textbf{engañosa} con clases desbalanceadas. Si el $95\%$ de los datos son clase negativa, un clasificador que predice siempre negativo tiene accuracy del 95\% pero no detecta ningún positivo. En esos casos se prefieren $F_1$, AUC-ROC o AUC-PR.
\end{alerta}

\subsection{Curva ROC y AUC}

La \textbf{curva ROC} (\textit{Receiver Operating Characteristic}) grafica la Sensibilidad (TPR) contra la Tasa de Falsos Positivos (FPR $= 1 - \text{TNR}$) para distintos umbrales de clasificación $\tau \in [0,1]$.

El \textbf{AUC} (\textit{Area Under the Curve}) resume la curva en un único valor:
\begin{itemize}
  \item $\text{AUC} = 1$: clasificador perfecto.
  \item $\text{AUC} = 0.5$: clasificador aleatorio (sin poder discriminante).
  \item $\text{AUC} < 0.5$: peor que aleatorio.
\end{itemize}

\begin{conceptoclave}
  AUC-ROC mide la probabilidad de que el modelo asigne un score mayor a una observación positiva aleatoria que a una negativa aleatoria. Es invariante al umbral de decisión y al desbalance de clases.
\end{conceptoclave}

\subsection{Curva Precisión-Recall (PR)}

Con datos muy desbalanceados, la curva ROC puede ser optimista. La \textbf{curva PR} grafica Precisión vs.\ Recall para distintos umbrales; el \textbf{AP} (\textit{Average Precision}) es el área bajo ella. Es más informativa cuando los positivos son muy escasos.

\subsection{Extensión a multiclase}

Para $K > 2$ clases:
\begin{itemize}
  \item \textbf{Macro promedio:} promedio no ponderado de la métrica por clase.
  \item \textbf{Weighted promedio:} promedio ponderado por la frecuencia de cada clase.
  \item \textbf{Micro promedio:} agrega VP, FP, FN de todas las clases antes de calcular.
\end{itemize}

% ─────────────────────────────────────────────────────────────────────────────
\section{Clases Desbalanceadas}
% ─────────────────────────────────────────────────────────────────────────────

En muchos problemas reales (detección de fraude, diagnóstico médico, fallas de hardware) una clase es mucho más frecuente que la otra. Esto perjudica a la mayoría de los clasificadores, que optimizan accuracy global.

\subsection{Estrategias de resampling}

\subsubsection{Oversampling: SMOTE}

\textbf{SMOTE} (\textit{Synthetic Minority Oversampling Technique}, Chawla et al.\ 2002) genera observaciones sintéticas de la clase minoritaria interpolando entre observaciones reales:
\begin{enumerate}
  \item Para cada observación $\bx_i$ de la clase minoritaria, encontrar sus $k$ vecinos más cercanos.
  \item Generar un punto sintético $\tilde{\bx} = \bx_i + \lambda (\bx_{nn} - \bx_i)$ donde $\bx_{nn}$ es un vecino elegido al azar y $\lambda \sim \mathcal{U}(0,1)$.
\end{enumerate}

\subsubsection{Undersampling}

Reducir la clase mayoritaria eliminando observaciones al azar o mediante métodos más sofisticados (Tomek links, ENN).

\begin{alerta}
  SMOTE debe aplicarse \textbf{solo sobre el conjunto de entrenamiento}, nunca sobre el de validación o prueba. Aplicarlo antes de dividir los datos produce data leakage.
\end{alerta}

\subsection{Ajuste de pesos de clase}

Muchos clasificadores (RF, SVM, regresión logística) aceptan el parámetro \texttt{class\_weight} que pondera más los errores en la clase minoritaria durante el entrenamiento:
\[
  w_k = \frac{n}{K \cdot n_k}
\]
donde $n_k$ es el número de observaciones de la clase $k$.

% ─────────────────────────────────────────────────────────────────────────────
\section{Ingeniería de Features}
% ─────────────────────────────────────────────────────────────────────────────

La \textbf{ingeniería de features} (\textit{feature engineering}) es el proceso de transformar los datos crudos en representaciones más informativas para el modelo.

\subsection{Transformaciones habituales}

\begin{itemize}
  \item \textbf{Estandarización:} $x' = (x - \bar{x})/s$. Obligatoria para SVM y KNN; recomendable para redes neuronales.
  \item \textbf{Normalización min-max:} $x' = (x - x_{\min})/(x_{\max} - x_{\min})$.
  \item \textbf{Transformación logarítmica:} reduce asimetría en variables con larga cola derecha.
  \item \textbf{Codificación de variables categóricas:} one-hot encoding, label encoding, target encoding.
  \item \textbf{Variables de interacción:} $x_1 \cdot x_2$, capturan relaciones no lineales entre pares.
  \item \textbf{Binning (discretización):} convierte variables continuas en intervalos.
\end{itemize}

\subsection{Selección de features}

\begin{itemize}
  \item \textbf{Filtrado:} basado en estadísticos univariados (correlación, $\chi^2$, información mutua).
  \item \textbf{Wrapper:} selección basada en el rendimiento del modelo (RFE -- \textit{Recursive Feature Elimination}).
  \item \textbf{Embedded:} regularización LASSO, importancia en Random Forest.
\end{itemize}

% ─────────────────────────────────────────────────────────────────────────────
\section{Comparación de Modelos: Grid Search y Pipelines}
% ─────────────────────────────────────────────────────────────────────────────

\subsection{Búsqueda de hiperparámetros}

Los hiperparámetros no se aprenden durante el entrenamiento; deben fijarse externamente. Las estrategias de búsqueda son:

\begin{itemize}
  \item \textbf{Grid search:} evalúa todas las combinaciones de una grilla prefijada de valores.
  \item \textbf{Random search:} muestrea combinaciones al azar de distribuciones de valores; más eficiente que grid search para espacios grandes.
  \item \textbf{Búsqueda bayesiana:} utiliza un modelo probabilístico del rendimiento para guiar la búsqueda.
\end{itemize}

\begin{observacion}
  La evaluación durante la búsqueda de hiperparámetros debe realizarse con validación cruzada \textbf{sobre el conjunto de entrenamiento}. El conjunto de prueba queda reservado para la evaluación final del modelo seleccionado.
\end{observacion}

\subsection{Pipelines}

Un \textbf{pipeline} encadena pasos de preprocesamiento y modelado en un único objeto, garantizando que:
\begin{enumerate}
  \item Las transformaciones se aprenden solo sobre el fold de entrenamiento en cada iteración de CV.
  \item No hay filtración de información del conjunto de validación.
  \item El código de producción es reproducible y no presenta discrepancias respecto al entrenamiento.
\end{enumerate}

\begin{conceptoclave}
  El uso de pipelines es una práctica esencial de ML responsable: evita el \textbf{data leakage} y simplifica el despliegue del modelo.
\end{conceptoclave}

% ─────────────────────────────────────────────────────────────────────────────
\section*{Resumen de la Unidad}
\addcontentsline{toc}{section}{Resumen de la Unidad}
% ─────────────────────────────────────────────────────────────────────────────

\begin{center}
\begin{tabular}{p{3cm} p{4cm} p{3.5cm} p{3.5cm}}
  \toprule
  \textbf{Modelo} & \textbf{Supuesto clave} & \textbf{Ventaja} & \textbf{Limitación} \\
  \midrule
  Árbol de decisión & Sin supuestos distribucionales & Interpretable & Alta varianza \\
  Random Forest & IID, suficientes features & Robusto, importancia & Caja negra \\
  SVM & Separabilidad (kernel) & Eficaz en alta dim. & Costoso en $n$ grande \\
  KNN & Proximidad implica similitud & Sin entrenamiento & Maldición dim. \\
  MLP & Aproximador universal & Muy flexible & Requiere muchos datos \\
  \bottomrule
\end{tabular}
\end{center}

% ─────────────────────────────────────────────────────────────────────────────
\section*{Referencias}
\addcontentsline{toc}{section}{Referencias}
% ─────────────────────────────────────────────────────────────────────────────

\begin{itemize}
  \item Breiman, L. (2001). \textit{Random Forests}. Machine Learning, 45(1), 5--32.
  \item Chawla, N.\ V., et al.\ (2002). \textit{SMOTE: Synthetic Minority Over-sampling Technique}. JAIR, 16, 321--357.
  \item Géron, A. (2022). \textit{Hands-On Machine Learning with Scikit-Learn, Keras, and TensorFlow} (3rd ed.). O'Reilly Media.
  \item Hastie, T., Tibshirani, R., \& Friedman, J. (2009). \textit{The Elements of Statistical Learning} (2nd ed.). Springer.
  \item James, G., et al.\ (2021). \textit{An Introduction to Statistical Learning with Applications in R} (2nd ed.). Springer.
\end{itemize}

% ─────────────────────────────────────────────────────────────────────────────
\end{document}
